\documentclass[conference]{IEEEtran}
\usepackage{amsmath,amssymb}
\usepackage{booktabs}
\usepackage{url}
\usepackage[hidelinks]{hyperref}

\begin{document}

\title{CinemaAI: Design and Offline Evaluation of a Sentence-Embedding
Content-Based and User-Based Collaborative Movie Recommender for a Cinema
Ticketing Platform}

\author{\IEEEauthorblockN{Author Name}
\IEEEauthorblockA{Institution\\
City, Country\\
email@example.com}}

\maketitle

\begin{abstract}
Cinema ticketing platforms increasingly rely on recommendation systems to
surface relevant films to their audiences. This paper presents CinemaAI, a
production movie recommendation service deployed as a FastAPI microservice
within a three-tier cinema management system, and reports the first systematic
offline evaluation of its algorithms. The service combines (i) a content-based
recommender that encodes movie metadata with the \texttt{all-MiniLM-L6-v2}
sentence-embedding model and ranks by cosine similarity, (ii) a user-based
collaborative filtering (CF) model with $k=20$ nearest neighbors and
similarity-weighted rating prediction, and (iii) a two-tier popularity fallback
for cold-start users. We evaluate all three models under an identical per-user
temporal 80/10/10 split on two datasets: the MovieLens~1M benchmark (1{,}000{,}209
ratings, 6{,}040 users) and a fully synthetic CinemaAI dataset (55{,}775
interactions, 600 users) used for pipeline validation. Performance is measured
with Precision@10, Recall@10, NDCG@10, MAP@10, catalog coverage, RMSE, and MAE.
The results reveal three actionable findings. First, the value of content-based
recommendation depends strongly on metadata richness: with full metadata
(overview, cast, keywords) the sentence-embedding recommender outperforms all
alternatives by 3--4$\times$, whereas with title-and-genre metadata only
(MovieLens~1M) it degrades to near the CF level. Second, the production CF
formula---which ranks by raw similarity-weighted predicted rating without
mean-centering or support weighting---underperforms even an item-mean baseline
in rating prediction (RMSE 1.081 vs.\ 0.993 on MovieLens~1M). Third, under a
temporal split on real data, the popularity baseline is remarkably strong for
ranking (Precision@10 $=0.039$) but recommends only 3\% of the catalog,
justifying its role as a fallback rather than a primary model. We discuss
concrete, formula-level improvements and release the evaluation harness
alongside the system.
\end{abstract}

\begin{IEEEkeywords}
recommender systems, collaborative filtering, content-based filtering,
sentence embeddings, offline evaluation, MovieLens
\end{IEEEkeywords}

\section{Introduction}
Online cinema ticketing platforms present users with a catalog of currently
showing and upcoming films. Unlike large streaming catalogs, a cinema catalog
is small and volatile---titles enter and leave weekly---yet the recommendation
problem remains: a visitor deciding what to watch benefits from ``similar
movies'' on a detail page and from a personalized ``recommended for you'' shelf
on the home page.

CinemaAI is the recommendation and conversational-AI subsystem of CinePremier,
a full-stack cinema management system consisting of a React frontend, a Spring
Boot backend, and a Python FastAPI microservice. The microservice implements
two complementary recommenders that mirror the classical dichotomy in the
literature: a content-based model over movie metadata and a user-based
collaborative filtering model over explicit ratings, with a popularity fallback
for cold-start users.

While the system is functional in production, its algorithms had never been
evaluated quantitatively. This paper closes that gap. Our contributions are:
\begin{enumerate}
  \item \textbf{A documented production architecture} for embedding-based movie
  recommendation in a microservice setting, including offline embedding
  precomputation and graceful degradation (Section~\ref{sec:system}).
  \item \textbf{A reproducible offline evaluation harness} implementing
  Precision@K, Recall@K, NDCG@K, MAP@K, catalog coverage, RMSE, and MAE under a
  per-user temporal split, with the production formulas ported verbatim
  (Section~\ref{sec:setup}).
  \item \textbf{An empirical comparison} of the three production models on
  MovieLens~1M (primary, real-world benchmark) and a synthetic CinemaAI dataset
  (secondary, pipeline validation), yielding formula-level diagnoses of the
  production CF model and evidence that metadata richness governs content-based
  performance (Sections~\ref{sec:results}--\ref{sec:discussion}).
\end{enumerate}

\section{Related Work}
\textbf{Collaborative filtering.} Neighborhood-based CF dates to
GroupLens~\cite{resnick1994grouplens}, which predicted a user's rating as a
similarity-weighted average of neighbors' ratings. Sarwar et
al.~\cite{sarwar2001item} introduced the item-based variant; Herlocker et
al.~\cite{herlocker2004evaluating} established evaluation methodology for CF
systems, including the importance of mean-centering and significance (support)
weighting in neighborhood models---both absent from the CinemaAI production
formula, a gap our results quantify.

\textbf{Content-based filtering.} Classical content-based recommenders
represent items with TF-IDF bag-of-words vectors~\cite{ricci2015handbook}.
Modern systems replace sparse lexical vectors with dense semantic embeddings;
Sentence-BERT~\cite{reimers2019sbert} produces sentence-level embeddings
suitable for cosine-similarity ranking, and MiniLM~\cite{wang2020minilm}
provides a distilled 6-layer encoder (\texttt{all-MiniLM-L6-v2}, 384
dimensions) that trades little accuracy for large speedups, making it practical
for CPU-only production services such as ours.

\textbf{Evaluation.} Cremonesi et al.~\cite{cremonesi2010performance} showed
that top-$N$ accuracy and rating-prediction accuracy can disagree, and that
non-personalized popularity is a strong top-$N$ baseline---an effect we
reproduce under a temporal split. Harper and
Konstan~\cite{harper2015movielens} describe the MovieLens datasets used here.

\section{System Overview}\label{sec:system}
CinePremier is a three-tier system:
\begin{itemize}
  \item \textbf{Frontend (React/Vite).} The movie detail page displays a
  ``similar movies'' grid fed by the content-based endpoint; the home page
  displays a personalized ``recommended for you'' shelf fed by the
  collaborative endpoint, falling back to the public catalog when no
  personalized results exist.
  \item \textbf{Backend (Spring Boot).} Exposes
  \texttt{/api/v1/recommendation/content/\{movieId\}} and
  \texttt{/api/v1/recommendation/collaborative/\{userId\}}; a thin
  \texttt{AIClient} proxies to the AI microservice and returns an empty list on
  failure, so recommendation outages never break page rendering.
  \item \textbf{AI microservice (FastAPI, Python).} Hosts the two recommenders
  and a retrieval-augmented movie chatbot. Movie embeddings are precomputed
  offline by a batch script from the PostgreSQL catalog and stored in a pickle
  file; at query time only the query-side encoding and cosine ranking are
  performed.
\end{itemize}

\section{Recommendation Algorithms}\label{sec:algorithms}

\subsection{Content-Based Recommendation with Sentence Embeddings}
Each movie $m$ is represented by a text document concatenating its metadata:
$d_m = \text{description} \oplus \text{director} \oplus \text{actors} \oplus
\text{genres}$. The document is encoded by the pretrained sentence-embedding
model \texttt{all-MiniLM-L6-v2}~\cite{reimers2019sbert,wang2020minilm} into a
dense vector $\mathbf{v}_m \in \mathbb{R}^{384}$. Given a query movie $q$, all
catalog movies are ranked by cosine similarity
\begin{equation}
\cos(\mathbf{v}_q,\mathbf{v}_m)=
\frac{\mathbf{v}_q\cdot\mathbf{v}_m}{\lVert\mathbf{v}_q\rVert\,\lVert\mathbf{v}_m\rVert},
\end{equation}
the query movie is excluded, and the top-10 are returned. Unlike the TF-IDF
representation common in earlier systems, the embedding captures semantic
similarity between synonymous descriptions.

\subsection{User-Based Collaborative Filtering}
Each user $u$ is a sparse rating vector over the set $I_u$ of movies they have
rated. Similarity between two users is the cosine over their ratings (the dot
product ranges over co-rated movies only):
\begin{equation}
\mathrm{sim}(u,v)=
\frac{\sum_{m\in I_u\cap I_v} r_{u,m}\,r_{v,m}}
{\sqrt{\sum_{m\in I_u} r_{u,m}^2}\sqrt{\sum_{m\in I_v} r_{v,m}^2}}.
\end{equation}
For a target user $u$, the $k=20$ most similar users with positive similarity
form the neighborhood $N(u)$. The predicted rating for an unseen movie $m$ is
the similarity-weighted average
\begin{equation}\label{eq:cfpred}
\hat{r}_{u,m}=
\frac{\sum_{v\in N(u)}\mathrm{sim}(u,v)\,r_{v,m}}
{\sum_{v\in N(u)}\mathrm{sim}(u,v)},
\end{equation}
and movies already rated or booked by the user are excluded. The top-10 movies
by predicted rating are returned. Note that this production formula uses
\emph{raw} ratings (no mean-centering) and ranks purely by predicted value
regardless of how many neighbors support the prediction;
Section~\ref{sec:discussion} shows both choices are measurably harmful.

\subsection{Cold-Start Fallback}
Users without ratings receive a two-tier non-personalized list: (1) movies
ranked by mean review rating; if empty, (2) movies ranked by the number of
paid bookings (popularity).

\section{Datasets}\label{sec:datasets}
\textbf{MovieLens 1M (primary).} The main experiments use the MovieLens~1M
stable benchmark~\cite{harper2015movielens}: 1{,}000{,}209 anonymous ratings
(integers 1--5) of 3{,}883 movies from 6{,}040 users, each with at least 20
ratings. Interactions are ordered chronologically per user; the earliest 80\%
form the training set (797{,}758 interactions), the next 10\% validation, and
the latest 10\% the test set (102{,}759 interactions). This temporal strategy
prevents leakage of future interactions into training. Because MovieLens
provides only titles and genres as item metadata, the content-based document is
$d_m=\text{title}\oplus\text{genres}$.

\textbf{Synthetic CinemaAI dataset (secondary).} A separately generated
synthetic dataset (seed 20260717) is used only to validate metadata processing,
the evaluation pipeline, and application behavior under the \emph{full}
metadata schema of the production system (overview, director, cast, keywords,
genres): 1{,}200 movies, 600 users, 55{,}775 interactions with explicit ratings
(0.5--5.0) and a binary \texttt{liked} flag, split per-user temporally 80/10/10
(44{,}376 / 5{,}560 / 5{,}839). Every record carries
\texttt{synthetic\_flag=true}. Results from synthetic data are reported
separately and are \emph{not} interpreted as evidence of real-world
effectiveness.

\section{Experimental Setup}\label{sec:setup}
\textbf{Protocol.} All models are fit on the training split only. For each
user, the candidate set is the full catalog minus the user's training items
(mirroring the production seen-item filter); each model produces a top-10
ranking. A held-out test item is \emph{relevant} if its rating $\ge 4$ on
MovieLens~1M (standard practice), and if $\texttt{liked}=1$ (primary) or rating
$\ge 3.5$ (sensitivity variant) on the synthetic dataset. Users with no
relevant test items are excluded from ranking metrics.

\textbf{Models.} (1)~\emph{Popularity}: a single ranking by training
interaction count, personalized only by the seen-item filter.
(2)~\emph{User-based CF} ($k=20$): the production formula of
Eq.~\eqref{eq:cfpred}, ported verbatim (dictionary implementation for the
synthetic dataset; a mathematically equivalent vectorized implementation for
MovieLens~1M). (3)~\emph{Content-based SBERT}: because the production endpoint
is item-to-item, per-user evaluation requires an adaptation---each user is
represented by a profile vector equal to the rating-weighted mean of the
embeddings of their liked training items (falling back to all training items
when none are liked), and candidates are ranked by cosine to this profile.

\textbf{Metrics.} Precision@10, Recall@10, NDCG@10, and MAP@10 averaged over
users; catalog coverage@10 (fraction of the catalog recommended to at least one
user); RMSE and MAE for rating prediction, compared against global-mean and
per-item-mean baselines. All computations are deterministic; two independent
runs produce byte-identical results.

\textbf{Environment.} Python~3.11, sentence-transformers~3.0.1 (CPU),
numpy/pandas. Encoding 3{,}883 MovieLens documents takes 7\,s; the full
MovieLens evaluation completes in under one minute.

\section{Results}\label{sec:results}

\begin{table}[t]
\caption{Top-10 ranking quality, MovieLens 1M\\(relevance: rating $\ge 4$; 5{,}827 users)}
\label{tab:ml_ranking}
\centering
\begin{tabular}{lccccc}
\toprule
Model & P@10 & R@10 & NDCG & MAP & Cov. \\
\midrule
Popularity & \textbf{.0394} & \textbf{.0446} & \textbf{.0506} & \textbf{.0227} & .0301 \\
User CF $k{=}20$ & .0067 & .0096 & .0100 & .0043 & \textbf{.7005} \\
Content SBERT & .0060 & .0113 & .0093 & .0040 & .2856 \\
\bottomrule
\end{tabular}
\end{table}

\begin{table}[t]
\caption{Rating prediction, MovieLens 1M test set}
\label{tab:ml_rating}
\centering
\begin{tabular}{lcc}
\toprule
Model & RMSE & MAE \\
\midrule
Global mean & 1.1634 & 0.9575 \\
Item mean & \textbf{0.9925} & \textbf{0.7893} \\
User-based CF $k{=}20$ & 1.0807 & 0.8421 \\
\bottomrule
\end{tabular}
\end{table}

\begin{table}[t]
\caption{Top-10 ranking quality, synthetic dataset\\(relevance: liked $=1$; 363 users)}
\label{tab:syn_ranking}
\centering
\begin{tabular}{lccccc}
\toprule
Model & P@10 & R@10 & NDCG & MAP & Cov. \\
\midrule
Popularity & .0036 & .0195 & .0124 & .0084 & .0133 \\
User CF $k{=}20$ & .0039 & .0233 & .0121 & .0067 & \textbf{.5317} \\
Content SBERT & \textbf{.0127} & \textbf{.0787} & \textbf{.0417} & \textbf{.0245} & .3942 \\
\bottomrule
\end{tabular}
\end{table}

\begin{table}[t]
\caption{Rating prediction, synthetic test set}
\label{tab:syn_rating}
\centering
\begin{tabular}{lcc}
\toprule
Model & RMSE & MAE \\
\midrule
Global mean & 0.8404 & 0.6795 \\
Item mean & \textbf{0.8221} & \textbf{0.6641} \\
User-based CF $k{=}20$ & 0.9532 & 0.7543 \\
\bottomrule
\end{tabular}
\end{table}

\subsection{MovieLens 1M (Real-World Benchmark)}
Tables~\ref{tab:ml_ranking} and~\ref{tab:ml_rating} report ranking and
rating-prediction quality on MovieLens~1M. Popularity dominates all ranking
metrics; the two personalized models are statistically comparable to each other
and far behind. In rating prediction, CF improves on the global mean but
remains worse than the per-item mean.

\subsection{Synthetic CinemaAI Dataset (Pipeline Validation)}
Tables~\ref{tab:syn_ranking} and~\ref{tab:syn_rating} report the same
quantities on the synthetic dataset. The content-based model leads every
ranking metric by 3--4$\times$. The sensitivity variant (relevance: rating
$\ge 3.5$; 524 users) preserves the ordering: content-based leads with
P@10 $= 0.0116$ against $0.0053$ for both alternatives.

\section{Discussion}\label{sec:discussion}

\textbf{F1 --- Metadata richness governs content-based performance.} The
sentence-embedding recommender is the best model by 3--4$\times$ on the
synthetic dataset, where documents contain overview, cast, keywords, director,
and genres---the schema of the production system. On MovieLens~1M, where only
title and genres are available, the same model collapses to CF-level
performance (P@10 $= 0.0060$). The practical implication for CinemaAI is
direct: the investment in rich catalog metadata is precisely what makes the
embedding approach viable, and enriching metadata is likely a higher-leverage
improvement than changing the model.

\textbf{F2 --- The production CF formula underperforms trivial baselines.} On
both datasets, ranking by raw similarity-weighted predicted rating is no better
than (synthetic) or far worse than (MovieLens) the popularity baseline, and
rating prediction is worse than the per-item mean (RMSE 1.081 vs.\ 0.993 on
MovieLens; 0.953 vs.\ 0.822 on synthetic). Two textbook
defects~\cite{herlocker2004evaluating} explain this: (i)~no
\emph{mean-centering}---users with different rating scales are averaged
directly; the standard fix is
\begin{equation}
\hat{r}_{u,m}=\bar{r}_u+
\frac{\sum_{v}\mathrm{sim}(u,v)\,(r_{v,m}-\bar{r}_v)}
{\sum_{v}\mathrm{sim}(u,v)};
\end{equation}
(ii)~no \emph{support weighting}---an item rated 5 by a single neighbor
outranks an item rated 4.5 by fifteen neighbors, which also explains CF's very
high coverage (70\% of the catalog) being a symptom of noise rather than
healthy diversity. Both fixes are one-line formula changes in the production
service.

\textbf{F3 --- Popularity is a strong ranking baseline under temporal splits,
but a poor recommender.} On MovieLens~1M, popularity dominates all personalized
models on every ranking metric, consistent
with~\cite{cremonesi2010performance} and amplified by the temporal split
(recent test interactions concentrate on currently popular titles). Yet it
recommends only 3.0\% of the catalog---every user receives nearly the same
list. This supports its production role as a cold-start \emph{fallback} while
cautioning against reading its accuracy as user value.

\textbf{F4 --- Rating-prediction and ranking quality disagree.} CF beats the
global mean on RMSE yet ranks worse than popularity; content-based has no
rating semantics at all yet ranks best on the synthetic dataset. This
reproduces the classical observation of~\cite{cremonesi2010performance} that
RMSE-centric development can mislead top-$N$ system design, and justifies our
metric suite.

\textbf{Implications for the production system.} In order of expected
leverage: (1)~apply mean-centering and a minimum-support threshold (or
shrinkage) to the CF predictor; (2)~blend content and CF scores into a hybrid
ranker, with the blend weight tuned on the validation split; (3)~keep
popularity strictly as fallback; (4)~for larger catalogs, replace the $O(N)$
brute-force cosine scan with an approximate-nearest-neighbor index.

\section{Threats to Validity}
\emph{Construct.} Offline relevance (rating $\ge 4$, liked $=1$) is a proxy
for user satisfaction; online A/B testing remains necessary before claiming
user-facing impact. The content-based model required an item-to-item
$\rightarrow$ user-profile adaptation for per-user evaluation; production
behavior on the movie detail page is item-to-item and is not directly measured
here. \emph{Internal.} Hyperparameters were fixed at production values
($k=20$, $K=10$) rather than tuned on the validation split; tuned baselines
could shift absolute numbers, though the formula-level defects of F2 are
parameter-independent in direction. \emph{External.} The synthetic dataset is
generated, carries \texttt{synthetic\_flag=true}, and is used only to validate
the pipeline under the full metadata schema; no claim about real audience
behavior is derived from it. MovieLens users are movie-enthusiast raters and
may differ from cinema-ticket buyers. \emph{Reliability.} All experiments are
deterministic and were verified to reproduce exactly across independent runs;
the split, seed (20260717), and harness are fixed and released.

\section{Conclusion and Future Work}
We presented the CinemaAI recommendation service and its first systematic
offline evaluation. The study delivered a reusable evaluation harness and three
formula-level findings: sentence-embedding content-based recommendation is
highly effective exactly when catalog metadata is rich; the production CF
predictor needs mean-centering and support weighting before it can compete
with trivial baselines; and popularity, while a strong offline ranking baseline
under temporal evaluation, is unsuitable as a primary recommender due to
negligible coverage. Future work includes implementing the corrected CF
formula, learning a validation-tuned hybrid blend of content and collaborative
scores, replacing brute-force similarity search with an ANN index, and
validating the offline findings with an online experiment on the CinePremier
platform.

\section*{Acknowledgments}
MovieLens data courtesy of GroupLens Research. The synthetic dataset was
generated for engineering validation and is marked as such in every record.

\begin{thebibliography}{9}

\bibitem{ricci2015handbook}
F.~Ricci, L.~Rokach, and B.~Shapira, Eds., \emph{Recommender Systems Handbook},
2nd~ed. Springer, 2015.

\bibitem{harper2015movielens}
F.~M. Harper and J.~A. Konstan, ``The MovieLens datasets: History and
context,'' \emph{ACM Trans. Interact. Intell. Syst.}, vol.~5, no.~4,
Article~19, 2015.

\bibitem{reimers2019sbert}
N.~Reimers and I.~Gurevych, ``Sentence-BERT: Sentence embeddings using Siamese
BERT-networks,'' in \emph{Proc. EMNLP-IJCNLP}, 2019, pp. 3982--3992.

\bibitem{wang2020minilm}
W.~Wang, F.~Wei, L.~Dong, H.~Bao, N.~Yang, and M.~Zhou, ``MiniLM: Deep
self-attention distillation for task-agnostic compression of pre-trained
transformers,'' in \emph{Proc. NeurIPS}, 2020.

\bibitem{resnick1994grouplens}
P.~Resnick, N.~Iacovou, M.~Suchak, P.~Bergstrom, and J.~Riedl, ``GroupLens: An
open architecture for collaborative filtering of netnews,'' in \emph{Proc.
CSCW}, 1994, pp. 175--186.

\bibitem{sarwar2001item}
B.~Sarwar, G.~Karypis, J.~Konstan, and J.~Riedl, ``Item-based collaborative
filtering recommendation algorithms,'' in \emph{Proc. WWW}, 2001, pp.
285--295.

\bibitem{herlocker2004evaluating}
J.~L. Herlocker, J.~A. Konstan, L.~G. Terveen, and J.~T. Riedl, ``Evaluating
collaborative filtering recommender systems,'' \emph{ACM Trans. Inf. Syst.},
vol.~22, no.~1, pp. 5--53, 2004.

\bibitem{cremonesi2010performance}
P.~Cremonesi, Y.~Koren, and R.~Turrin, ``Performance of recommender algorithms
on top-N recommendation tasks,'' in \emph{Proc. RecSys}, 2010, pp. 39--46.

\end{thebibliography}

\end{document}
